% ---
% title: Two exams, conditionally
% tags: conditional-probability
% ---

\begin{problem}
	We continue Exercise~\exref{two-exams}. A randomly selected student of an
	Economics class passed the Math exam ($A$) with probability $0.72$, the
	Microeconomics exam ($B$) with probability $0.64$, and both exams with
	probability $0.54$.
	(a) The student passed Microeconomics. What is the probability that they
	also passed Math?
	(b) The student passed Math. What is the probability that they did not
	pass Microeconomics?
	(c) The student passed at least one exam. What is the probability that they
	passed Math?
	(d) The student passed at least one exam. What is the probability that they
	did not pass Microeconomics?
\end{problem}

\begin{slideproblem}
	As in Exercise~\exref{two-exams}: $P(A)=0.72$ (Math), $P(B)=0.64$
	(Micro), $P(AB)=0.54$. Probability that the student (a) passed Math, if
	they passed Micro; (b) failed Micro, if they passed Math; (c) passed Math,
	if they passed at least one exam; (d) failed Micro, if they passed at least
	one exam?
\end{slideproblem}

\begin{slides}
	\small
	\exstatement
	\begin{columns}[T]
		\begin{column}{.62\textwidth}
			\uncover<2->{\sol
				$P(E\mid F)=\dfrac{P(EF)}{P(F)}$: $F$ happened, so $F$ is the
				new sample space.}\\[1mm]
			\uncover<3->{(a) Given: passed Micro. New sample space $B$
				(shaded).}\\[.5mm]
			\uncover<4->{$P(A\mid B)=\dfrac{P(AB)}{P(B)}=\dfrac{0.54}{0.64}
					=\dfrac{27}{32}\approx 0.844$.}\\[1mm]
			\uncover<5->{(b) Given: passed Math. New sample space $A$.}\\[.5mm]
			\uncover<6->{$P(B^c\mid A)=\dfrac{P(AB^c)}{P(A)}=\dfrac{0.18}{0.72}
					=0.25$.}
		\end{column}
		\begin{column}{.36\textwidth}
			\centering
			\uncover<2->{%
				\begin{tikzpicture}[scale=.55]
					\only<3-4>{\fill[teal!25] (1.6,0) circle (1.5);}
					\only<5->{\fill[navy!20] (0,0) circle (1.5);}
					\draw[thick] (-2,-2) rectangle (3.6,2.1);
					\draw[navy,thick] (0,0) circle (1.5);
					\draw[teal,thick] (1.6,0) circle (1.5);
					\node at (-.65,0) {\scriptsize 0.18};
					\node at (.8,0) {\scriptsize 0.54};
					\node at (2.25,0) {\scriptsize 0.10};
					\node at (2.95,-1.6) {\scriptsize 0.18};
					\node[navy] at (-1.4,1.55) {$A$};
					\node[teal] at (3.0,1.55) {$B$};
					\node at (.8,1.75) {\scriptsize $\Omega$};
				\end{tikzpicture}}
		\end{column}
	\end{columns}
\end{slides}

\begin{slides}
	\small
	\begin{columns}[T]
		\begin{column}{.62\textwidth}
			(c) Given: passed at least one exam. New sample space
			$A\cup B$ (shaded):
			$P(A\cup B)=0.18+0.54+0.10=0.82$.\\[1.5mm]
			\uncover<2->{$A$ lies inside $A\cup B$, so $A(A\cup B)=A$:}
			\uncover<3->{\[ P(A\mid A\cup B)=\frac{P(A)}{P(A\cup B)}
					=\frac{0.72}{0.82}=\frac{36}{41}\approx 0.878. \]}%
			\uncover<4->{(d) Failed Micro \emph{and} passed at least one exam:
				only Math, $B^c(A\cup B)=AB^c$:}
			\uncover<5->{\[ P(B^c\mid A\cup B)=\frac{P(AB^c)}{P(A\cup B)}
					=\frac{0.18}{0.82}=\frac{9}{41}\approx 0.220. \]}
		\end{column}
		\begin{column}{.36\textwidth}
			\centering
			\begin{tikzpicture}[scale=.55]
				\fill[gold!30] (0,0) circle (1.5);
				\fill[gold!30] (1.6,0) circle (1.5);
				\draw[thick] (-2,-2) rectangle (3.6,2.1);
				\draw[navy,thick] (0,0) circle (1.5);
				\draw[teal,thick] (1.6,0) circle (1.5);
				\node at (-.65,0) {\scriptsize 0.18};
				\node at (.8,0) {\scriptsize 0.54};
				\node at (2.25,0) {\scriptsize 0.10};
				\node at (2.95,-1.6) {\scriptsize 0.18};
				\node[navy] at (-1.4,1.55) {$A$};
				\node[teal] at (3.0,1.55) {$B$};
				\node at (.8,1.75) {\scriptsize $\Omega$};
			\end{tikzpicture}
		\end{column}
	\end{columns}
	\uncover<6->{\smallskip
		\alert{Knowing more changes the probability:} $P(A)=0.72$, but
		$P(A\mid B)\approx0.844$ and $P(A\mid A\cup B)\approx0.878$.}
\end{slides}

\begin{solution}
	\idea Conditional probability:
	\[ P(E\mid F)=\frac{P(EF)}{P(F)}. \]
	We are told that $F$ happened, so only the outcomes in $F$ are still
	possible: $F$ becomes the new, smaller sample space. Inside it we measure
	how large the part where $E$ also happens is.

	\begin{minipage}{.56\linewidth}
		\step{The Venn diagram.} From Exercise~\exref{two-exams}
		($A$ = passed Math, $B$ = passed Microeconomics):
		\begin{itemize}[nosep]
			\item both exams: $P(AB)=0.54$;
			\item only Math: $P(AB^c)=0.72-0.54=0.18$;
			\item only Micro: $P(A^cB)=0.64-0.54=0.10$;
			\item at least one: $P(A\cup B)=0.72+0.64-0.54=0.82$;
			\item neither: $1-0.82=0.18$.
		\end{itemize}
	\end{minipage}\hfill
	\begin{minipage}{.4\linewidth}
		\centering
		\begin{tikzpicture}[scale=.7]
			\draw[thick] (-2,-2) rectangle (3.6,2.1);
			\draw[navy,thick] (0,0) circle (1.5);
			\draw[teal,thick] (1.6,0) circle (1.5);
			\node at (-.65,0) {\small 0.18};
			\node at (.8,0) {\small 0.54};
			\node at (2.25,0) {\small 0.10};
			\node at (2.95,-1.6) {\small 0.18};
			\node[navy] at (-1.4,1.55) {$A$};
			\node[teal] at (3.0,1.55) {$B$};
			\node at (.8,1.75) {\small $\Omega$};
		\end{tikzpicture}
	\end{minipage}

	\step{(a)} We know that the student passed Microeconomics, so the new
	sample space is $B$, the circle $B$ in the diagram, with probability
	$0.64$. Inside $B$, the student also passed Math in the region $AB$, with
	probability $0.54$. So
	\[ P(A\mid B)=\frac{P(AB)}{P(B)}=\frac{0.54}{0.64}=\frac{27}{32}=0.84375. \]

	\step{(b)} Now the condition is ``passed Math'': the new sample space is
	the circle $A$, with probability $0.72$. Inside $A$, the student did not
	pass Microeconomics in the region ``only Math'', $AB^c$, with probability
	$P(A)-P(AB)=0.72-0.54=0.18$. So
	\[ P(B^c\mid A)=\frac{P(AB^c)}{P(A)}=\frac{0.18}{0.72}=0.25. \]
	Equivalently, $P(B\mid A)=0.54/0.72=0.75$ and $P(B^c\mid A)=1-0.75=0.25$:
	for a fixed condition, the complement rule works as usual.

	\step{(c)} The condition is ``passed at least one exam'', the event
	$A\cup B$, with probability $0.82$: the two circles together. We need the
	part of $A\cup B$ where the student passed Math, that is
	$A\cap(A\cup B)$. Every student who passed Math passed at least one exam,
	so $A$ lies entirely inside $A\cup B$, and the intersection is just $A$:
	\[ P(A\mid A\cup B)=\frac{P(A\cap(A\cup B))}{P(A\cup B)}
		=\frac{P(A)}{P(A\cup B)}=\frac{0.72}{0.82}=\frac{36}{41}\approx 0.878. \]

	\step{(d)} Again the new sample space is $A\cup B$. Inside it, the student
	did not pass Microeconomics only in the region ``only Math'': passing at
	least one exam but not Micro means passing Math but not Micro. So
	$B^c\cap(A\cup B)=AB^c$, with probability $0.18$:
	\[ P(B^c\mid A\cup B)=\frac{P(AB^c)}{P(A\cup B)}=\frac{0.18}{0.82}
		=\frac{9}{41}\approx 0.220. \]
	\begin{answer}
		(a) $\frac{27}{32}=0.84375$ \quad (b) $0.25$ \quad
		(c) $\frac{36}{41}\approx 0.878$ \quad (d) $\frac{9}{41}\approx 0.220$
	\end{answer}
	\begin{remark}
		\textbf{Common mistake:} dividing by the wrong probability. The
		denominator is always the probability of the condition (what we
		\emph{know}), never of the event we ask about. In (a), $0.54/0.72=0.75$
		would be $P(B\mid A)$, the answer to a different question.
	\end{remark}
	\begin{remark}
		The information changes the probabilities: unconditionally the student
		passed Math with probability $0.72$, but knowing that they passed
		Microeconomics raises this to $0.844$. In (d), failing Microeconomics has
		probability $0.36$ unconditionally, but only $0.220$ among students who
		passed at least one exam.
	\end{remark}
\end{solution}
